% !TEX program = lualatex
\documentclass[aspectratio=43,professionalfonts,handout]{beamer}
\usepackage{beamer_template}

\newcommand{\LectureNum}{12}
\newcommand{\LectureTitle}{無理関数}
\newcommand{\TextPages}{p.\,94-95}
\newcommand{\NextTopic}{逆関数}
\newcommand{\NextPages}{p.\,96-97}
\newcommand{\CourseName}{基礎数学B}
\renewcommand{\TermName}{前期}

\begin{document}

% -----------------------------------------------
% スライド1: 表紙＋目標＋用語・定義
% -----------------------------------------------
\begin{frame}{\CourseName\quad 第\LectureNum 回「\LectureTitle 」}
  {\small \TermName\quad \TeacherName\quad ／\quad 教科書 \TextPages\quad
  \textbf{目標}：無理関数の定義域とグラフを理解し、変形グラフが描ける}

  \begin{block}{用語・定義}
  \begin{description}[labelwidth=5.5em]
    \item[\kw{無理関数}] 根号（$\sqrt{\phantom{x}}$）の中に変数を含む関数
    \item[\kw{定義域}] 根号内 $P(x)\geq 0$ を解いて求める（例：$y=\sqrt{x-1}$ は $x\geq 1$）
    \item[\kw{$y=\sqrt{x}$ のグラフ}] 放物線 $x=y^2$（$y\geq 0$）の半分、第1象限の曲線
  \end{description}
  \end{block}
\end{frame}

% -----------------------------------------------
% スライド2: 公式・性質
% -----------------------------------------------
\begin{frame}{公式・性質}
  \begin{block}{}
  \begin{itemize}
    \item $x$ 軸対称：$y=f(x)$ $\to$ $y=-f(x)$
    \item $y$ 軸対称：$y=f(x)$ $\to$ $y=f(-x)$
    \item 原点対称：$y=f(x)$ $\to$ $y=-f(-x)$
    \item 平行移動：$y=\sqrt{x-p}+q$（始点が $(p,\,q)$ に移動）
  \end{itemize}
  \end{block}
\end{frame}

% -----------------------------------------------
% スライド3: 例1→問1
% -----------------------------------------------
\begin{frame}{例1→問1}
  \begin{exampleblock}{例5) p.94（定義域を求めよ）}
    （板書で解説）
  \end{exampleblock}

  \vfill

  \begin{block}{問5) p.94（定義域3問）\quad 自分で解いてみよう}
    （ノートに解くこと）
  \end{block}
\end{frame}

% -----------------------------------------------
% スライド4: 例2→問2
% -----------------------------------------------
\begin{frame}{例2→問2}
  \begin{exampleblock}{例題2) p.95（グラフ・定義域・値域）／ 例) p.95（対称移動）}
    （板書で解説）
  \end{exampleblock}

  \vfill

  \begin{block}{問6) p.95（グラフ3問）、問7・問8) p.95（対称移動）\quad 自分で解いてみよう}
    （ノートに解くこと）
  \end{block}
\end{frame}

% -----------------------------------------------
% まとめ
% -----------------------------------------------
\begin{frame}{まとめ}
  \begin{itemize}
    \item $y=\sqrt{P(x)}$ の定義域は $P(x)\geq 0$ を解いて求める
    \item $y=\sqrt{x}$ のグラフは放物線 $x=y^2$ の上半分
    \item $x$ 軸・$y$ 軸・原点に関する対称移動の公式を使い分ける
  \end{itemize}

  \vfill
  {\small 基本問題：175, 176, 177, 178}
\end{frame}

\end{document}
